% ID: SEL-D0068-OA-E01
% ESTADO: validada. S068.
\textbf{a)} Multiplicamos la segunda inecuación por $15$. Las cotas quedan
$(15-3y)/5\le x\le12-3y$, con $y\ge1$. Al comparar las cotas resulta $y\le15/4$. En ese intervalo la cota inferior de $x$ es positiva. Para $y=1$ obtenemos $(12/5,1)$ y $(9,1)$. Las dos rectas oblicuas se cortan en $(3/4,15/4)$, el tercer vértice.
\begin{center}
\setlength{\unitlength}{0.9mm}
\begin{picture}(98,60)(-8,-6)
\qbezier(0,0)(40,0)(80,0)
\qbezier(0,0)(0,24)(0,48)
\qbezier(19.2,10)(45.6,10)(72,10)
\qbezier(72,10)(39,23.75)(6,37.5)
\qbezier(6,37.5)(12.6,23.75)(19.2,10)
\put(81,0){$x$}\put(2,48){$y$}
\put(15,4){\scriptsize $(12/5,1)$}
\put(71,4){\scriptsize $(9,1)$}
\put(8,39){\scriptsize $(3/4,15/4)$}
\end{picture}
\end{center}
\textbf{b)} Como $F=5(x+3y)$, alcanza su máximo $\boxed{60}$ en todo el lado $x+3y=12$, entre $(9,1)$ y $(3/4,15/4)$. En el vértice restante toma el valor $5(12/5+3)=27$. Al ser una función lineal, el mínimo es $\boxed{27}$ y se alcanza sólo en $(12/5,1)$.
